\includegraphics{Lab6/Group12/Lab6/traces/L6-2-6.throughput.png}\\

De grafiek toont een sterke overeenkomst tussen de simulatie (groene punten) en de analytische curve $T = G \cdot e^{-2G}$ (blauwe lijn).\\
Vorm van de curves: beide tonen hetzelfde gedrag: T stijgt aanvankelijk met G, bereikt een maximum, en daalt daarna. Dit bevestigt dat bij lage load meer stations zorgen voor meer throughput, maar bij hoge load zijn er veel collisions en daalt de efficiëntie sterk.\\

Maximum van de gesimuleerde throughput: de simulatiepunten bereiken een maximum van ongeveer $T \approx 0.175 – 0.184$ rond $G \approx 0.4 – 0.5$, wat goed overeenkomt met het theoretische maximum van:
$T_{max} = \dfrac{1}{2e} \approx 0.184$ bij $G = 0.5$